
\renewcommand{\fundamentacion}{El Espacio Curricular Dibujo Técnico tiene por finalidad en el Ciclo Básico o Primer Ciclo promover en los estudiantes capacidades, conocimientos, habilidades, destrezas, valores y actitudes propias de este espacio curricular en función del perfil profesional técnico. En este espacio los estudiantes adquieren los fundamentos del dibujo técnico para poder interpretar y elaborar información gráfica, y aplicar en forma correcta los sistemas de representación para relacionar el espacio con el plano. Teniendo en cuenta que el Dibujo Técnico es de desarrollo eminentemente práctico, las propuestas de enseñanza aprendizaje deberán poner a los estudiantes en situación de aplicación de los saberes construidos. Los docentes deben generar diferentes actividades formativas para que los estudiantes expresen las soluciones gráficas, con claridad, precisión y objetividad. Además, potenciar en los estudiantes el uso de los instrumentos propios de la disciplina con habilidad y destreza, como medio de transmisión de las ideas científico-técnicas. Las propuestas de enseñanza aprendizaje deberán facilitar el trabajo autónomo de los estudiantes, potenciar las técnicas de indagación e investigación, y asegurar las aplicaciones y transferencias de lo aprendido a la vida real. El Dibujo Técnico se utiliza como una herramienta en distintos espacios curriculares de la Educación Técnica, por ello es fundamental que en el Ciclo Básico ó Primer Ciclo los estudiantes construyan saberes en el trazado y croquizado a mano alzada. Los procesos de aprendizaje por lo tanto, están en función del “saber hacer” de saberes propios del dibujo técnico, y de su integración con los saberes del resto de los espacios curriculares. Para ello, el docente deberá utilizar estrategias tales como la realización de croquis acotados en base a modelos reales y la identificación de elementos normalizados en planos técnicos ya ejecutados, entre otros.}

\renewcommmand{\materia}{Dibujo Técnico}
\renewcommand{\escuela}{9-239 Cerro el Sosneado}
